\documentclass[output=paper]{langscibook}
\ChapterDOI{10.5281/zenodo.22880290}
\title{\textit{Seige} – Fishing with a net fixed on two mangrove poles}
\author{Enoch Horai Magum}
\abstract{}
\begin{document}
\maketitle\label{chap:31}
\begin{figure}[H]
    \includegraphics[trim=0 30 0 20,clip,width=.89\linewidth]{figures/31_1.jpg}
    \caption{\textit{Seige} -- Fishing with a net fixed on two mangrove poles}
    \label{fig:31_1}
\end{figure}
\begin{xltabular}{\textwidth}{lXX}
	(\ref{ex:text24-1})&\textit{O meho kaku vagana vaa tea kave en teara, a ta Teapu, na dao raara bono seige.}&Another kind of fishing with nets here with us, the Teop people, is called \textit{seige}.\\
	(\ref{ex:text24-2})&\textit{A vagana bona saka vaapenapena haana bea peha gea buaku ge bero aba nao tea nao vagana.}&This kind of fishing allows one or two or many people to go fishing.\\
	(\ref{ex:text24-3})&\textit{O kaku vagana bona he na paku vonaen:}&This kind of fishing is done as follows:\\
	(\ref{ex:text24-4})&\textit{Eara na vaahae rara a kave, eara repaa hua vo tea maa hum to vavaasun kave komana raara.}&We take the net on board, then we paddle to the places where we put up nets.\\
	(\ref{ex:text24-5})&\textit{Amaa hum bona he toro hum vai to gunagunaha kaakoo rio iana beve tarahi batana.}&These places must be places where the fish go when the tide is ebbing.\\
	(\ref{ex:text24-6})&\textit{O kaku hum bona he eiei tea maa hoa.}&This kind of place is in the passages.\\
	(\ref{ex:text24-7})&\textit{Beara suguna, eara repaa busu nia kave, eara repaa rahi ge vaasun pina nie bona hanana teo iana.}&When we arrive, we throw the net, then we pull it or put it up so that it blocks the way of the fish.\\
	(\ref{ex:text24-8})&\textit{Eara repaa pam nio buo punete na kave teo buo boon to paa vaasun vuara, eara he tea mei tea vaagunaha a kave.}&Then we tie the two ends of the net to two mangrove poles that we have erected before we let down the net.\\
	(\ref{ex:text24-9})&\textit{Beve kavara, eara repaa vahakaa pina koa maara o iana to mene gunaha kaakoo teo meho hum.}&After that we beat the water to prevent the fish from going down to another place.\\
	(\ref{ex:text24-10})&\textit{Bea tahii taneo bata mana tea tarahi, eove o vuri to kahi taneo hana rio iana.}&When the sea eventually starts ebbing, this is the time when the fish start hanging in the net.\\
	(\ref{ex:text24-11})&\textit{Beori hana rori, eara na antee batara tea no ropusu kahi riori bona kave, ge be hiki, ei mana, beve tarahi vakavara.}&When they are hanging there, we can take them out of the net, or if they don’t, this eventually means that it is low tide.\\
	(\ref{ex:text24-12})&\textit{Evehee, eara toro usu bata koa ni rara o iana vai o bebeera to mene buno kurukurus bona kave, a kave repaa taparata.}&But we must watch the big fish that they do not move so much that the net gets opened.\\
	(\ref{ex:text24-13})&\textit{O kaku vagana vai o mataa tabae eara sa hua binubinu haara tea rake ge vatevatete a iana.}&It is a good kind of fishing, because we do not paddle around to search for or chase the fish.\\
	(\ref{ex:text24-14})&\textit{Eara kahi sun ihuana koa nom beori hana rori tea kave.}&We would just stand still waiting for them to hang in the net.\\
	(\ref{ex:text24-15})&\textit{Bea tahii masi vakavara, eara repaa voosu vahaa vaan.}&When the sea is at low tide, we go home back to the village.\\
\end{xltabular}
\end{document}